% Part: many-valued-logic
% Chapter: infinite-valued-logics
% Section: lukasiewicz

\documentclass[../../../include/open-logic-section]{subfiles}

\begin{document}

\olfileid{mvl}{inf}{luk}

\olsection{Logika \L ukasiewicz}

\begin{editorial}
  Iki mung draf ringkes saka perangan bab logika
  \L ukasiewicz mawa nilai tanpa wates.
\end{editorial}

\begin{defn}\ollabel{def:lukasiewicz} Logika \L ukasiewicz mawa nilai tanpa wates
~$\LogLuk[\infty]$ ditetepake nganggo matriks iki:
\begin{enumerate}
  \item Basa proposisional baku $\Lang L_0$ kang ngemot
  $\lfalse$, $\lnot$, $\land$, $\lor$, $\lif$.
  \item Himpunan nilai kabeneran $V_\infty$.
  \item $1$ mung siji-sijine nilai pinilih, yaiku $V^+ = \{1\}$.
  \item Kanggo tetapan salah, $\tf{\lfalse}[\LogLuk[\infty]] = 0$. Fungsi kabeneran liyane diwenehake ing ngisor iki:
  \begin{align*}
    \tf{\lnot}[\LogLuk](x) & = 1 - x\\
    \tf{\land}[\LogLuk](x,y) & = \min(x,y)\\
    \tf{\lor}[\LogLuk](x,y) & = \max(x,y)\\
    \tf{\lif}[\LogLuk](x,y) & = \min(1,1-(x-y)) = \begin{cases}
      1 & \text{yen } x \le y\\
      1-(x-y) & \text{yen ora.}
    \end{cases}
    \end{align*}
\end{enumerate}
Logika \L ukasiewicz $m$ nilai ditetepake kanthi cara padha, nanging $V = V_m$.
\end{defn}

\begin{prop}
  Logika $\LogLuk[3]$ kang ditetepake ing
  \olref[thr][luk]{def:lukasiewicz} padha karo $\LogLuk[3]$
  kang ditetepake ing \olref{def:lukasiewicz}.
\end{prop}

\begin{proof}
  Nilai tetapan salah uga cocog ing loro matriks. Iki katon
  kanthi mbandhingake tabel kabeneran konektif ing
  \olref[thr][luk]{def:lukasiewicz} karo tabel kang ditemtokake
  dening persamaan ing \olref{def:lukasiewicz}:
  \begin{center}
    \begin{tabular}{c|c} 
      $\tf{\lnot}$ & \\ 
      \hline  
      $1$ & $0$ \\ 
      $1/2$ & $1/2$ \\
      $0$ & $1$ 
    \end{tabular}
    \quad
    \begin{tabular}{c|ccc} 
      $\tf{\land}[\LogLuk[3]]$ & $1$ & $1/2$ & $0$ \\ 
      \hline 
      $1$ & $1$ & $1/2$ & $0$ \\ 
      $1/2$ & $1/2$ & $1/2$ & $0$\\ 
      $0$ & $0$ & $0$ & $0$ 
    \end{tabular}
    \\[2ex]
    \begin{tabular}{c|ccc} 
      $\tf{\lor}[\LogLuk[3]]$ & $1$ & $1/2$ & $0$ \\ 
      \hline 
      $1$ & $1$ & $1$ & $1$ \\ 
      $1/2$ & $1$ & $1/2$ & $1/2$ \\
      $0$ & $1$ & $1/2$ & $0$ 
    \end{tabular}
    \quad
    \begin{tabular}{c|ccc} 
      $\tf{\lif}[\LogLuk[3]]$ & $1$ & $1/2$ & $0$ \\ 
      \hline 
      $1$ & $1$ & $1/2$ & $0$ \\ 
      $1/2$ & $1$ & $1$ & $1/2$  \\ 
      $0$ & $1$ & $1$ & $1$ 
    \end{tabular}
  \end{center} 
\end{proof}

\begin{prop}\ollabel{prop:luk-infty-m}
  Yen $\Gamma \Entails[\LogLuk[\infty]] !B$, mula $\Gamma
  \Entails[\LogLuk[m]] !B$ kanggo kabeh~$m \ge 2$.
\end{prop}

\begin{proof}
  Latihan.
\end{proof}

\begin{prob}
  Buktèkna \olref[mvl][inf][luk]{prop:luk-infty-m}.
\end{prob}

Satemene arah kosok baline uga bener.

Logika \L ukasiewicz mawa nilai tanpa wates iku logika fuzzy
kang paling kawentar. Ing pustaka logika fuzzy, kondisional asring
ditetepake minangka $\lnot !A \lor !B$. Asile yaiku logika
Kleene kang kuwat kanthi nilai tanpa wates.

\begin{prob}
  Tuduhna yen $(p \lif q) \lor (q \lif p)$ iku tautologi
  ing~$\LogLuk[\infty]$.
\end{prob}

\end{document}
